\documentclass[12pt,a4paper]{report}

% Packages
\usepackage[utf8]{inputenc}
\usepackage[T1]{fontenc}
\usepackage{geometry}
\usepackage{graphicx}
\usepackage{hyperref}
\usepackage{xcolor}
\usepackage{listings}
\usepackage{booktabs}
\usepackage{array}
\usepackage{longtable}
\usepackage{enumitem}
\usepackage{fancyhdr}
\usepackage{titlesec}
\usepackage{tocloft}
\usepackage{amsmath}
\usepackage{microtype}
\usepackage{setspace}

% Page geometry
\geometry{
  top=2.5cm,
  bottom=2.5cm,
  left=3cm,
  right=2.5cm
}

% Colors
\definecolor{primaryblue}{RGB}{30, 80, 162}
\definecolor{codegreen}{RGB}{40, 120, 60}
\definecolor{codegray}{RGB}{128, 128, 128}
\definecolor{codebackground}{RGB}{245, 245, 250}
\definecolor{mongored}{RGB}{0, 104, 74}
\definecolor{reactblue}{RGB}{97, 218, 251}

% Hyperref config
\hypersetup{
  colorlinks=true,
  linkcolor=primaryblue,
  urlcolor=primaryblue,
  citecolor=primaryblue,
  pdftitle={MERN School Management System - Project Report},
  pdfauthor={Rohith Uppunuthula},
  pdfsubject={Full Stack Web Application Documentation}
}

% Code listing style
\lstdefinestyle{codestyle}{
  backgroundcolor=\color{codebackground},
  commentstyle=\color{codegreen},
  keywordstyle=\color{primaryblue}\bfseries,
  numberstyle=\tiny\color{codegray},
  stringstyle=\color{mongored},
  basicstyle=\ttfamily\small,
  breakatwhitespace=false,
  breaklines=true,
  captionpos=b,
  keepspaces=true,
  numbers=left,
  numbersep=5pt,
  showspaces=false,
  showstringspaces=false,
  showtabs=false,
  tabsize=2,
  frame=single,
  rulecolor=\color{codegray},
  xleftmargin=10pt,
  xrightmargin=5pt
}
\lstset{style=codestyle}

% Section formatting
\titleformat{\chapter}[display]
  {\normalfont\huge\bfseries\color{primaryblue}}
  {\chaptertitlename\ \thechapter}{20pt}{\Huge}

\titleformat{\section}
  {\normalfont\Large\bfseries\color{primaryblue}}
  {\thesection}{1em}{}

\titleformat{\subsection}
  {\normalfont\large\bfseries\color{primaryblue!80}}
  {\thesubsection}{1em}{}

% Header/Footer
\pagestyle{fancy}
\fancyhf{}
\rhead{\textcolor{primaryblue}{\small MERN School Management System}}
\lhead{\textcolor{codegray}{\small Project Report}}
\rfoot{\thepage}
\lfoot{\textcolor{codegray}{\small Rohith Uppunuthula}}
\renewcommand{\headrulewidth}{0.4pt}
\renewcommand{\footrulewidth}{0.4pt}

% Line spacing
\onehalfspacing

%=============================================================================
\begin{document}

%-----------------------------------------------------------------------------
% Title Page
%-----------------------------------------------------------------------------
\begin{titlepage}
  \centering
  \vspace*{2cm}

  {\Huge\bfseries\color{primaryblue} MERN School Management System}\\[0.6cm]
  {\large\color{codegray} Full Stack Web Application}\\[0.3cm]
  \rule{\linewidth}{0.5pt}\\[1.5cm]

  {\large\bfseries Project Report}\\[2cm]

  \begin{tabular}{ll}
    \textbf{Author}     & Rohith Uppunuthula \\[4pt]
    \textbf{Technology} & MERN Stack (MongoDB, Express.js, React.js, Node.js) \\[4pt]
    \textbf{Frontend}   & React 18, Vite, Redux Toolkit, Material UI \\[4pt]
    \textbf{Backend}    & Node.js, Express.js, Mongoose \\[4pt]
    \textbf{Database}   & MongoDB (Local / Atlas) \\[4pt]
    \textbf{Date}       & April 2026 \\[4pt]
    \textbf{Repository} & \href{https://github.com/rohith-uppunuthula/MERN-School-Management-System}{GitHub} \\
  \end{tabular}

  \vfill

  \rule{\linewidth}{0.5pt}\\[0.5cm]
  {\small\color{codegray}
    Streamline school management, class organization, and add students and faculty.\\
    Seamlessly track attendance, assess performance, and provide feedback.
  }
\end{titlepage}

%-----------------------------------------------------------------------------
% Table of Contents
%-----------------------------------------------------------------------------
\tableofcontents
\newpage

%=============================================================================
\chapter{Project Overview}
%=============================================================================

\section{Introduction}

The \textbf{MERN School Management System} is a comprehensive, role-based web application designed to digitize and streamline day-to-day school operations. Built on the industry-standard MERN stack, it unifies the workflows of three distinct user types --- \textit{Administrators}, \textit{Teachers}, and \textit{Students} --- into a single, cohesive platform accessible from any modern web browser.

The system replaces manual record-keeping with automated, real-time data management: attendance is recorded and visualized per-subject; examination results are entered and available instantly; complaints and notices flow through structured channels; and every user has a personalized dashboard tailored to their role.

\section{Objectives}

\begin{itemize}[leftmargin=2em]
  \item Provide a centralized portal for school administrators to manage classes, subjects, teachers, and students.
  \item Enable teachers to record attendance per subject, assess student performance, and communicate via complaints.
  \item Allow students to monitor their own attendance, subject-wise marks, and receive official notices.
  \item Deliver interactive data visualizations (bar charts, pie charts) to support academic decision-making.
  \item Ensure clean separation of concerns through role-based routing and Redux state management.
\end{itemize}

\section{Key Features at a Glance}

\begin{longtable}{@{}p{4.5cm}p{10cm}@{}}
  \toprule
  \textbf{Feature} & \textbf{Description} \\
  \midrule
  \endhead
  Role-Based Access & Three distinct roles: Admin, Teacher, Student — each with its own dashboard and permissions. \\[4pt]
  Admin Dashboard & Add/remove students, teachers, classes, and subjects; view all data. \\[4pt]
  Attendance Tracking & Teachers mark presence/absence per subject per date; data stored per student. \\[4pt]
  Exam Result Management & Teachers enter marks per subject; students view their own results. \\[4pt]
  Notice Board & Admin creates announcements visible to all users in the school. \\[4pt]
  Complaints System & Teachers and students submit complaints; admin can review them. \\[4pt]
  Data Visualization & Interactive pie and bar charts for attendance and performance data. \\[4pt]
  Guest Access & Visitors can explore the app as a guest without registration. \\[4pt]
  Responsive UI & Material UI components adapt to desktop and mobile viewports. \\[4pt]
  Secure Auth & bcrypt password hashing; role-aware login flows. \\
  \bottomrule
\end{longtable}

%=============================================================================
\chapter{System Architecture}
%=============================================================================

\section{High-Level Architecture}

The application follows a classic three-tier architecture:

\begin{enumerate}[leftmargin=2em]
  \item \textbf{Presentation Tier} --- React.js SPA built with Vite, served from \texttt{localhost:4823} in development. Redux Toolkit manages global application state; React Router v6 handles client-side navigation.
  \item \textbf{Application Tier} --- Node.js + Express.js REST API served from \texttt{localhost:5000}. Controllers handle business logic; Mongoose ODM abstracts database operations.
  \item \textbf{Data Tier} --- MongoDB document database (local \texttt{mongod} or Atlas cloud). Six collections map to the six Mongoose schemas.
\end{enumerate}

\begin{figure}[h]
  \centering
  \begin{tabular}{|c|}
    \hline
    \textbf{Browser (React SPA)} \\
    Vite dev server / Static hosting \\
    \hline
    $\downarrow$ HTTP / REST API $\uparrow$ \\
    \hline
    \textbf{Express.js API Server} \\
    Node.js runtime, port 5000 \\
    \hline
    $\downarrow$ Mongoose ODM $\uparrow$ \\
    \hline
    \textbf{MongoDB Database} \\
    Local or Atlas Cloud \\
    \hline
  \end{tabular}
  \caption{Three-Tier Architecture Overview}
\end{figure}

\section{Frontend Architecture}

\subsection{Technology Stack}

\begin{itemize}[leftmargin=2em]
  \item \textbf{React 18} --- Component-based UI library
  \item \textbf{Vite} --- Next-generation build tool and dev server (port 4823)
  \item \textbf{Redux Toolkit} --- Global state management with slices for each entity
  \item \textbf{React Router v6} --- Declarative client-side routing with role-based guards
  \item \textbf{Material UI (MUI)} --- Pre-built, accessible component library
  \item \textbf{Recharts} --- Composable charting library for data visualization
\end{itemize}

\subsection{Directory Structure}

\begin{lstlisting}[language=bash, caption=Frontend source tree]
frontend/src/
├── App.jsx                  # Root component with role-based routing
├── index.jsx                # React DOM entry point
├── index.css                # Global styles
├── assets/                  # Static images and SVGs
├── components/              # Shared UI components
│   ├── CustomBarChart.jsx
│   ├── CustomPieChart.jsx
│   ├── TableTemplate.jsx
│   ├── TableViewTemplate.jsx
│   ├── SeeNotice.jsx
│   ├── Popup.jsx
│   ├── ErrorPage.jsx
│   ├── AccountMenu.jsx
│   └── SpeedDialTemplate.jsx
├── pages/
│   ├── Homepage.jsx         # Public landing page
│   ├── LoginPage.jsx        # Shared login (all roles)
│   ├── ChooseUser.jsx       # Role selector
│   ├── admin/               # Admin-specific pages & dashboards
│   ├── teacher/             # Teacher-specific pages & dashboards
│   └── student/             # Student-specific pages & dashboards
└── redux/
    ├── store.jsx             # Redux store configuration
    ├── userRelated/          # Auth slice + thunks
    ├── studentRelated/       # Student slice + thunks
    ├── teacherRelated/       # Teacher slice + thunks
    ├── sclassRelated/        # Class slice + thunks
    ├── noticeRelated/        # Notice slice + thunks
    └── complainRelated/      # Complain slice + thunks
\end{lstlisting}

\subsection{Redux State Slices}

Each domain entity has its own Redux slice following the same pattern: a slice file defines the state shape and reducers; a handle file exports async thunks that call the REST API.

\begin{longtable}{@{}p{4.5cm}p{9.5cm}@{}}
  \toprule
  \textbf{Slice} & \textbf{Responsibility} \\
  \midrule
  \endhead
  \texttt{userSlice} & Authentication state, current role, current user object \\[4pt]
  \texttt{studentSlice} & Student list, student detail, loading/error states \\[4pt]
  \texttt{teacherSlice} & Teacher list, teacher detail \\[4pt]
  \texttt{sclassSlice} & Class list, class detail, students in class \\[4pt]
  \texttt{noticeSlice} & Notice list for current school \\[4pt]
  \texttt{complainSlice} & Complaint list for current school \\
  \bottomrule
\end{longtable}

\section{Backend Architecture}

\subsection{Technology Stack}

\begin{itemize}[leftmargin=2em]
  \item \textbf{Node.js} --- JavaScript runtime
  \item \textbf{Express.js 4} --- Minimal and flexible web framework
  \item \textbf{Mongoose 7} --- MongoDB object modelling for Node.js
  \item \textbf{bcrypt 5} --- Password hashing library (salt rounds based)
  \item \textbf{dotenv} --- Environment variable management
  \item \textbf{cors} --- Cross-Origin Resource Sharing middleware
  \item \textbf{nodemon} --- Auto-restart on file changes during development
\end{itemize}

\subsection{Directory Structure}

\begin{lstlisting}[language=bash, caption=Backend source tree]
backend/
├── index.js                 # App entry point (Express + Mongoose setup)
├── package.json
├── .env                     # MONGO_URL, SECRET_KEY
├── routes/
│   └── route.js             # All API route definitions
├── controllers/
│   ├── admin-controller.js
│   ├── student_controller.js
│   ├── teacher-controller.js
│   ├── class-controller.js
│   ├── subject-controller.js
│   ├── notice-controller.js
│   └── complain-controller.js
└── models/
    ├── adminSchema.js
    ├── studentSchema.js
    ├── teacherSchema.js
    ├── sclassSchema.js
    ├── subjectSchema.js
    ├── noticeSchema.js
    └── complainSchema.js
\end{lstlisting}

%=============================================================================
\chapter{Data Models}
%=============================================================================

\section{MongoDB Collections}

The database contains six collections, each managed by a Mongoose schema. Relationships are expressed via \texttt{ObjectId} references (similar to foreign keys in relational databases).

\section{Student Schema}

\begin{lstlisting}[language=JavaScript, caption=studentSchema.js]
{
  name:       String (required),
  rollNum:    Number (required),
  password:   String (required),
  sclassName: ObjectId -> 'sclass' (required),
  school:     ObjectId -> 'admin'  (required),
  role:       String  (default: "Student"),
  examResult: [
    {
      subName:       ObjectId -> 'subject',
      marksObtained: Number (default: 0)
    }
  ],
  attendance: [
    {
      date:    Date   (required),
      status:  Enum['Present', 'Absent'] (required),
      subName: ObjectId -> 'subject' (required)
    }
  ]
}
\end{lstlisting}

\section{Teacher Schema}

The teacher document links to a single subject and a single school (admin), enabling the system to know which teacher is responsible for which subject in which institution.

\begin{lstlisting}[language=JavaScript, caption=teacherSchema.js (summarized)]
{
  name:         String (required),
  email:        String (unique, required),
  password:     String (required),
  role:         String (default: "Teacher"),
  school:       ObjectId -> 'admin',
  teachSubject: ObjectId -> 'subject',
  teachSclass:  ObjectId -> 'sclass',
  attendance:   [{ date, presentCount, totalCount }]
}
\end{lstlisting}

\section{Admin Schema}

\begin{lstlisting}[language=JavaScript, caption=adminSchema.js (summarized)]
{
  name:       String (required),
  email:      String (unique, required),
  password:   String (required),
  role:       String (default: "Admin"),
  schoolName: String (required)
}
\end{lstlisting}

\section{Other Schemas Summary}

\begin{longtable}{@{}p{3.5cm}p{11cm}@{}}
  \toprule
  \textbf{Collection} & \textbf{Key Fields} \\
  \midrule
  \endhead
  \texttt{sclass} & \texttt{sclassName} (String, required), \texttt{school} (ObjectId $\to$ admin) \\[4pt]
  \texttt{subject} & \texttt{subName}, \texttt{sessions}, \texttt{sclassName} (ObjectId), \texttt{school} (ObjectId), \texttt{teacher} (ObjectId, optional) \\[4pt]
  \texttt{notice} & \texttt{title}, \texttt{details}, \texttt{date}, \texttt{school} (ObjectId) \\[4pt]
  \texttt{complain} & \texttt{user} (ObjectId), \texttt{date}, \texttt{complaint} (String), \texttt{school} (ObjectId) \\
  \bottomrule
\end{longtable}

\section{Entity Relationship Overview}

\begin{itemize}[leftmargin=2em]
  \item One \textbf{Admin} (school) owns many \textbf{Classes}, \textbf{Teachers}, \textbf{Students}, \textbf{Subjects}, \textbf{Notices}, and \textbf{Complaints}.
  \item Each \textbf{Class} belongs to one Admin and has many \textbf{Students} and \textbf{Subjects}.
  \item Each \textbf{Teacher} is assigned one \textbf{Subject} and one \textbf{Class}.
  \item Each \textbf{Student} belongs to one \textbf{Class} and has an embedded array of attendance records and exam results, both linked to subjects.
\end{itemize}

%=============================================================================
\chapter{API Reference}
%=============================================================================

\section{Authentication Endpoints}

\begin{longtable}{@{}p{1.5cm}p{5cm}p{7.5cm}@{}}
  \toprule
  \textbf{Method} & \textbf{Path} & \textbf{Description} \\
  \midrule
  \endhead
  POST & \texttt{/AdminReg} & Register a new school admin \\[3pt]
  POST & \texttt{/AdminLogin} & Admin authentication \\[3pt]
  GET  & \texttt{/Admin/:id} & Get admin/school details \\[3pt]
  POST & \texttt{/StudentReg} & Register a new student \\[3pt]
  POST & \texttt{/StudentLogin} & Student authentication \\[3pt]
  POST & \texttt{/TeacherReg} & Register a new teacher \\[3pt]
  POST & \texttt{/TeacherLogin} & Teacher authentication \\
  \bottomrule
\end{longtable}

\section{Student Endpoints}

\begin{longtable}{@{}p{1.5cm}p{6.5cm}p{6cm}@{}}
  \toprule
  \textbf{Method} & \textbf{Path} & \textbf{Description} \\
  \midrule
  \endhead
  GET    & \texttt{/Students/:id} & List all students of a school \\[3pt]
  GET    & \texttt{/Student/:id} & Get single student detail \\[3pt]
  PUT    & \texttt{/Student/:id} & Update student profile \\[3pt]
  DELETE & \texttt{/Student/:id} & Delete a single student \\[3pt]
  DELETE & \texttt{/Students/:id} & Delete all students of a school \\[3pt]
  DELETE & \texttt{/StudentsClass/:id} & Delete all students in a class \\[3pt]
  PUT    & \texttt{/StudentAttendance/:id} & Record attendance entry \\[3pt]
  PUT    & \texttt{/UpdateExamResult/:id} & Add/update exam marks \\[3pt]
  PUT    & \texttt{/RemoveAllStudentsSubAtten/:id} & Clear subject attendance (all) \\[3pt]
  PUT    & \texttt{/RemoveAllStudentsAtten/:id} & Clear all attendance (all) \\[3pt]
  PUT    & \texttt{/RemoveStudentSubAtten/:id} & Clear subject attendance (one) \\[3pt]
  PUT    & \texttt{/RemoveStudentAtten/:id} & Clear all attendance (one) \\
  \bottomrule
\end{longtable}

\section{Teacher Endpoints}

\begin{longtable}{@{}p{1.5cm}p{6.5cm}p{6cm}@{}}
  \toprule
  \textbf{Method} & \textbf{Path} & \textbf{Description} \\
  \midrule
  \endhead
  GET    & \texttt{/Teachers/:id} & List all teachers of a school \\[3pt]
  GET    & \texttt{/Teacher/:id} & Get single teacher detail \\[3pt]
  DELETE & \texttt{/Teacher/:id} & Delete a single teacher \\[3pt]
  DELETE & \texttt{/Teachers/:id} & Delete all teachers of a school \\[3pt]
  DELETE & \texttt{/TeachersClass/:id} & Delete all teachers of a class \\[3pt]
  PUT    & \texttt{/TeacherSubject} & Assign subject to teacher \\[3pt]
  POST   & \texttt{/TeacherAttendance/:id} & Record teacher attendance \\
  \bottomrule
\end{longtable}

\section{Class, Subject, Notice, and Complain Endpoints}

\begin{longtable}{@{}p{1.5cm}p{6.5cm}p{6cm}@{}}
  \toprule
  \textbf{Method} & \textbf{Path} & \textbf{Description} \\
  \midrule
  \endhead
  POST   & \texttt{/SclassCreate} & Create a new class \\[3pt]
  GET    & \texttt{/SclassList/:id} & List classes in a school \\[3pt]
  GET    & \texttt{/Sclass/:id} & Get class details \\[3pt]
  GET    & \texttt{/Sclass/Students/:id} & Get students in a class \\[3pt]
  DELETE & \texttt{/Sclass/:id} & Delete a class \\[3pt]
  DELETE & \texttt{/Sclasses/:id} & Delete all classes \\[3pt]
  \midrule
  POST   & \texttt{/SubjectCreate} & Create a subject \\[3pt]
  GET    & \texttt{/ClassSubjects/:id} & Subjects for a class \\[3pt]
  GET    & \texttt{/AllSubjects/:id} & All subjects in school \\[3pt]
  GET    & \texttt{/FreeSubjectList/:id} & Subjects without a teacher \\[3pt]
  GET    & \texttt{/Subject/:id} & Get subject detail \\[3pt]
  DELETE & \texttt{/Subject/:id} & Delete a subject \\[3pt]
  \midrule
  POST   & \texttt{/NoticeCreate} & Create a notice \\[3pt]
  GET    & \texttt{/NoticeList/:id} & List notices for school \\[3pt]
  PUT    & \texttt{/Notice/:id} & Update a notice \\[3pt]
  DELETE & \texttt{/Notice/:id} & Delete a single notice \\[3pt]
  DELETE & \texttt{/Notices/:id} & Delete all notices \\[3pt]
  \midrule
  POST   & \texttt{/ComplainCreate} & Submit a complaint \\[3pt]
  GET    & \texttt{/ComplainList/:id} & List complaints for school \\
  \bottomrule
\end{longtable}

%=============================================================================
\chapter{User Roles and Workflows}
%=============================================================================

\section{Role-Based Routing}

The \texttt{App.jsx} root component implements role-based routing using the Redux \texttt{currentRole} state. When the role is \texttt{null} (unauthenticated), public routes are shown. Upon successful login, the role is set and the appropriate dashboard is rendered exclusively.

\begin{lstlisting}[language=JavaScript, caption=App.jsx role-based routing logic]
const { currentRole } = useSelector(state => state.user);

// Public routes
{currentRole === null && <Routes>...</Routes>}

// Role-gated dashboards
{currentRole === "Admin"   && <AdminDashboard />}
{currentRole === "Student" && <StudentDashboard />}
{currentRole === "Teacher" && <TeacherDashboard />}
\end{lstlisting}

\section{Administrator Workflow}

The Admin is the system owner. After registering a school, an admin can:

\begin{enumerate}[leftmargin=2em]
  \item \textbf{Create Classes} --- e.g., ``Class 10A'', ``Class 12B''
  \item \textbf{Create Subjects} --- linked to a specific class (e.g., Mathematics for Class 10A)
  \item \textbf{Add Teachers} --- register teacher accounts and assign a subject and class
  \item \textbf{Add Students} --- register students into a class with a roll number
  \item \textbf{View All Records} --- students, teachers, classes, subjects, complaints
  \item \textbf{Manage Notices} --- create, update, and delete announcements
  \item \textbf{View Complaints} --- read all complaints submitted in the school
\end{enumerate}

\textbf{Admin Dashboard pages:}
\begin{itemize}[leftmargin=2em]
  \item \texttt{AdminHomePage} --- summary statistics and quick links
  \item \texttt{ShowClasses / AddClass / ClassDetails}
  \item \texttt{ShowStudents / AddStudent / ViewStudent / StudentAttendance / StudentExamMarks}
  \item \texttt{ShowTeachers / AddTeacher / TeacherDetails / ChooseClass / ChooseSubject}
  \item \texttt{ShowSubjects / SubjectForm / ViewSubject}
  \item \texttt{ShowNotices / AddNotice}
  \item \texttt{SeeComplains}
  \item \texttt{AdminProfile}
\end{itemize}

\section{Teacher Workflow}

\begin{enumerate}[leftmargin=2em]
  \item Log in using email and password
  \item View the assigned class and subject from the home page
  \item Navigate to \textbf{Class Details} to see enrolled students
  \item Take \textbf{Attendance} --- mark each student present or absent for the current date
  \item Enter \textbf{Exam Marks} per student per subject
  \item View a specific \textbf{Student Profile} with performance history
  \item Submit a \textbf{Complaint} if needed
\end{enumerate}

\section{Student Workflow}

\begin{enumerate}[leftmargin=2em]
  \item Log in using name, roll number, and password
  \item View the \textbf{Home Page} with a summary of attendance and performance
  \item Check \textbf{Subject Attendance} per subject with a pie chart visualization
  \item View \textbf{Exam Results} (marks per subject, bar chart)
  \item Read \textbf{Notices} from the school admin
  \item Submit a \textbf{Complaint}
  \item View and edit their \textbf{Profile}
\end{enumerate}

\section{Guest Access}

Unauthenticated visitors can choose ``Login as Guest'' from \texttt{ChooseUser}, which creates a guest session allowing exploration of the UI without any real data.

%=============================================================================
\chapter{Installation and Setup}
%=============================================================================

\section{Prerequisites}

\begin{itemize}[leftmargin=2em]
  \item Node.js $\geq$ 16.x
  \item npm $\geq$ 8.x
  \item MongoDB (local Community Server or Atlas account)
  \item Git
\end{itemize}

\section{Clone the Repository}

\begin{lstlisting}[language=bash, caption=Clone the project]
git clone https://github.com/rohith-uppunuthula/MERN-School-Management-System.git
cd MERN-School-Management-System
\end{lstlisting}

\section{Backend Setup}

\begin{lstlisting}[language=bash, caption=Backend installation and configuration]
cd backend
npm install

# Create environment file
cat > .env << 'EOF'
MONGO_URL=mongodb://127.0.0.1/smsproject
SECRET_KEY=secret123key
EOF

# Start the backend server (nodemon watches for changes)
npm start
# Server starts at http://localhost:5000
\end{lstlisting}

\section{Frontend Setup}

\begin{lstlisting}[language=bash, caption=Frontend installation and configuration]
cd frontend
npm install

# Create environment file
cat > .env << 'EOF'
REACT_APP_BASE_URL=http://localhost:5000
EOF

# Start Vite dev server
npm run dev
# App available at http://localhost:4823
\end{lstlisting}

\section{Database Configuration}

\subsection{Option 1: Local MongoDB}

\begin{lstlisting}[language=bash, caption=Start local MongoDB]
# Start mongod service
mongod

# Use connection string in .env
MONGO_URL=mongodb://127.0.0.1:27017/smsproject
\end{lstlisting}

\subsection{Option 2: MongoDB Atlas (Cloud)}

\begin{enumerate}[leftmargin=2em]
  \item Create an account at \href{https://mongodb.com/atlas}{mongodb.com/atlas}
  \item Create a free M0 cluster
  \item Click \textbf{Connect} $\to$ \textbf{Connect your application}
  \item Copy the connection string and replace in \texttt{.env}:
\end{enumerate}

\begin{lstlisting}[language=bash, caption=Atlas connection string]
MONGO_URL=mongodb+srv://<user>:<password>@<cluster>.mongodb.net/smsproject?retryWrites=true&w=majority
\end{lstlisting}

%=============================================================================
\chapter{Deployment}
%=============================================================================

\section{Backend Deployment on Render}

\begin{enumerate}[leftmargin=2em]
  \item Push the repository to GitHub.
  \item Create a new \textbf{Web Service} on \href{https://render.com}{Render}.
  \item Set the \textbf{Root Directory} to \texttt{backend}.
  \item Set \textbf{Build Command}: \texttt{npm install}
  \item Set \textbf{Start Command}: \texttt{npm start}
  \item Add environment variables: \texttt{MONGO\_URL} and \texttt{SECRET\_KEY}
  \item Render auto-deploys on every push to the linked branch.
\end{enumerate}

\section{Frontend Deployment on Netlify}

\begin{enumerate}[leftmargin=2em]
  \item Create a new \textbf{Netlify} project linked to the GitHub repository.
  \item Set \textbf{Base Directory}: \texttt{frontend}
  \item Set \textbf{Build Command}: \texttt{npm run build}
  \item Set \textbf{Publish Directory}: \texttt{dist} (Vite output)
  \item Add environment variable: \texttt{REACT\_APP\_BASE\_URL=https://your-backend.onrender.com}
  \item Netlify auto-builds on every push.
\end{enumerate}

\section{Frontend Deployment on Vercel}

\begin{enumerate}[leftmargin=2em]
  \item Import the GitHub repository into Vercel.
  \item Set \textbf{Root Directory}: \texttt{frontend}
  \item Same build command and output directory as Netlify.
  \item Set the same \texttt{REACT\_APP\_BASE\_URL} environment variable.
\end{enumerate}

%=============================================================================
\chapter{Security Considerations}
%=============================================================================

\section{Password Hashing}

All passwords are hashed using the \textbf{bcrypt} library before persistence to MongoDB. Plain-text passwords are never stored.

\begin{lstlisting}[language=JavaScript, caption=bcrypt usage in admin-controller.js]
const bcrypt = require('bcrypt');

// Registration
const salt = await bcrypt.genSalt(10);
const hashedPassword = await bcrypt.hash(password, salt);

// Login verification
const isMatch = await bcrypt.compare(plainPassword, storedHash);
\end{lstlisting}

\section{Environment Variables}

Sensitive configuration values (\texttt{MONGO\_URL}, \texttt{SECRET\_KEY}) are stored in \texttt{.env} files which are excluded from version control via \texttt{.gitignore}. The \texttt{dotenv} package loads these at runtime.

\section{CORS}

The Express server uses the \texttt{cors} middleware. In production, the CORS origin should be restricted to the deployed frontend URL to prevent unauthorized cross-origin requests.

\section{Known Security Improvements}

\begin{itemize}[leftmargin=2em]
  \item \textbf{JWT Authentication}: The \texttt{SECRET\_KEY} is present in the environment but full JWT middleware (token generation, verification, and route guarding) should be implemented to prevent unauthorized API access.
  \item \textbf{Input Validation}: Server-side schema validation (e.g., using \textit{Joi} or \textit{express-validator}) should be added to all POST/PUT routes.
  \item \textbf{Rate Limiting}: Add \texttt{express-rate-limit} to authentication endpoints to prevent brute-force attacks.
  \item \textbf{HTTPS}: Always serve the API over HTTPS in production (Render and Vercel provide this automatically).
\end{itemize}

%=============================================================================
\chapter{Dependencies}
%=============================================================================

\section{Backend Dependencies}

\begin{longtable}{@{}p{4cm}p{2.5cm}p{7.5cm}@{}}
  \toprule
  \textbf{Package} & \textbf{Version} & \textbf{Purpose} \\
  \midrule
  \endhead
  \texttt{express} & \^{}4.18.2 & Web framework for building the REST API \\[3pt]
  \texttt{mongoose} & \^{}7.0.4 & MongoDB ODM for schema definition and queries \\[3pt]
  \texttt{bcrypt} & \^{}5.1.0 & Secure password hashing \\[3pt]
  \texttt{cors} & \^{}2.8.5 & Cross-Origin Resource Sharing middleware \\[3pt]
  \texttt{dotenv} & \^{}16.0.3 & Load \texttt{.env} variables at runtime \\[3pt]
  \texttt{nodemon} & \^{}2.0.22 & Auto-restart server on file changes (dev) \\[3pt]
  \texttt{body-parser} & \^{}1.20.2 & Parse incoming request bodies (legacy, used as fallback) \\
  \bottomrule
\end{longtable}

\section{Frontend Key Dependencies}

\begin{longtable}{@{}p{5cm}p{9cm}@{}}
  \toprule
  \textbf{Package} & \textbf{Purpose} \\
  \midrule
  \endhead
  \texttt{react} / \texttt{react-dom} & UI library and DOM renderer \\[3pt]
  \texttt{vite} & Build tool and dev server \\[3pt]
  \texttt{@vitejs/plugin-react} & React fast-refresh plugin for Vite \\[3pt]
  \texttt{react-router-dom} & Client-side routing \\[3pt]
  \texttt{@reduxjs/toolkit} & Redux state management with slices \\[3pt]
  \texttt{react-redux} & React bindings for Redux \\[3pt]
  \texttt{@mui/material} & Material Design component library \\[3pt]
  \texttt{@mui/icons-material} & Material Design icon set \\[3pt]
  \texttt{recharts} & Composable charting library (pie, bar charts) \\[3pt]
  \texttt{axios} & HTTP client for API calls \\
  \bottomrule
\end{longtable}

%=============================================================================
\chapter{Project Structure Summary}
%=============================================================================

\section{Complete File Tree}

\begin{lstlisting}[language=bash, caption=Full project structure]
MERN-School-Management-System/
├── README.md
├── package-lock.json
├── backend/
│   ├── index.js               # Express app entry point
│   ├── package.json
│   ├── .env                   # MONGO_URL, SECRET_KEY
│   ├── seed.js                # (optional) seed data script
│   ├── routes/
│   │   └── route.js           # All REST API routes
│   ├── controllers/
│   │   ├── admin-controller.js
│   │   ├── student_controller.js
│   │   ├── teacher-controller.js
│   │   ├── class-controller.js
│   │   ├── subject-controller.js
│   │   ├── notice-controller.js
│   │   └── complain-controller.js
│   └── models/
│       ├── adminSchema.js
│       ├── studentSchema.js
│       ├── teacherSchema.js
│       ├── sclassSchema.js
│       ├── subjectSchema.js
│       ├── noticeSchema.js
│       └── complainSchema.js
└── frontend/
    ├── vite.config.js          # Vite configuration (port 4823)
    ├── package.json
    ├── .env                    # REACT_APP_BASE_URL
    └── src/
        ├── App.jsx             # Root routing component
        ├── index.jsx           # ReactDOM.createRoot entry
        ├── index.css           # Global styles
        ├── assets/             # Images, SVGs
        ├── components/         # Reusable shared components
        ├── pages/
        │   ├── Homepage.jsx
        │   ├── LoginPage.jsx
        │   ├── ChooseUser.jsx
        │   ├── Logout.jsx
        │   ├── admin/          # 15 admin pages
        │   ├── teacher/        # 7 teacher pages
        │   └── student/        # 7 student pages
        └── redux/
            ├── store.jsx
            ├── userRelated/
            ├── studentRelated/
            ├── teacherRelated/
            ├── sclassRelated/
            ├── noticeRelated/
            └── complainRelated/
\end{lstlisting}

%=============================================================================
\chapter{Conclusion}
%=============================================================================

\section{Summary}

The MERN School Management System demonstrates a production-ready full-stack JavaScript application that successfully addresses the core challenges of managing a school's digital operations. The project showcases:

\begin{itemize}[leftmargin=2em]
  \item \textbf{Clean separation of concerns} --- frontend state management via Redux, backend business logic in Express controllers, and data persistence in MongoDB schemas are kept strictly independent.
  \item \textbf{Role-based access control} --- three distinct user types with tailored views and restricted API access.
  \item \textbf{Rich data visualization} --- Recharts integration provides students and teachers with at-a-glance performance insights.
  \item \textbf{Developer experience} --- Vite's HMR, nodemon's auto-restart, and a clear project structure make the codebase easy to extend.
  \item \textbf{Flexible deployment} --- the decoupled architecture allows the backend and frontend to be hosted independently on Render, Netlify, Vercel, or any cloud provider.
\end{itemize}

\section{Potential Future Enhancements}

\begin{enumerate}[leftmargin=2em]
  \item \textbf{JWT Middleware} --- implement token-based authentication guards on all protected routes.
  \item \textbf{Real-time Notifications} --- integrate Socket.io for live notice and attendance updates.
  \item \textbf{File Uploads} --- allow profile pictures and document attachments via Cloudinary or S3.
  \item \textbf{Timetable Management} --- add a scheduling module for class timetables.
  \item \textbf{Parent Portal} --- extend role support to include parents who can view their child's progress.
  \item \textbf{Mobile App} --- port the frontend to React Native for iOS and Android.
  \item \textbf{Automated Testing} --- add Jest + React Testing Library unit tests and Playwright E2E tests.
  \item \textbf{CI/CD Pipeline} --- add GitHub Actions for automated linting, testing, and deployment.
\end{enumerate}

\vfill
\begin{center}
  \rule{0.5\linewidth}{0.4pt}\\[0.3cm]
  {\small\color{codegray}
    MERN School Management System --- Project Report\\
    Author: Rohith Uppunuthula $\cdot$ April 2026
  }
\end{center}

\end{document}
